The intuition that better pre-training data produces more generalizable language models drives substantial investment in corpus curation --- but this assumption has rarely been tested under controlled conditions. We evaluate OLMo-7B (Dolma, multi-stage curation) and Pythia-6.9B (The Pile, minimal curation) at matched training scale ($\sim$300B tokens), specifically testing whether the curated-corpus model achieves a higher MMLU/HellaSwag generalization balance ratio. Contrary to the hypothesis, Pythia achieves the higher ratio (0.565 vs.\ 0.538) --- a large, statistically robust reversal (95\% CI entirely against OLMo). The most informative finding is structural: both models achieve effectively identical HellaSwag commonsense scores at this training scale (0.458 each under 500-sample fast evaluation; full evaluation recommended to confirm), collapsing the ratio into a proxy for MMLU differences that are themselves architecture-confounded. This is consistent with a scale-dependent saturation of HellaSwag for 6--8B models at $\sim$300B tokens --- a hypothesis supported by the denominator's near-zero sensitivity to deduplication observed in prior work. We characterize the minimum requirements for future corpus quality studies: architecture-matched designs or temporal trajectory analysis to bound architecture confounds, and denominator-sensitivity verification before using ratio metrics as quality proxies.
